\documentclass[11pt,a4paper]{report}
\usepackage[italian]{babel}
\usepackage{fontspec}
\usepackage[margin=2.3cm,top=2.4cm,bottom=2.4cm]{geometry}
\usepackage{amsmath,amssymb,mathtools}
\usepackage{booktabs,array,tabularx,enumitem,multicol,xcolor}
\usepackage{tikz}
\usetikzlibrary{arrows.meta,positioning,calc,shapes.geometric}
\usepackage[most]{tcolorbox}
\usepackage{fancyhdr}
\usepackage{titlesec}
\usepackage[hidelinks,bookmarksnumbered]{hyperref}

\definecolor{cBlu}{HTML}{1F4E79}
\definecolor{cRosso}{HTML}{B03A2E}
\definecolor{cVerde}{HTML}{1E7B4B}
\definecolor{cGrigio}{HTML}{F2F2F2}
\definecolor{cArancio}{HTML}{B9770E}

\setlist{itemsep=2pt,topsep=3pt,parsep=0pt}
\setlength{\parindent}{0pt}
\setlength{\parskip}{5pt}
\renewcommand{\arraystretch}{1.18}

\titleformat{\chapter}[hang]{\Large\bfseries\color{cBlu}}{\thechapter.}{0.6em}{}
\titlespacing*{\chapter}{0pt}{-10pt}{14pt}
\titleformat{\section}{\large\bfseries\color{cBlu}}{\thesection}{0.6em}{}
\titleformat{\subsection}{\normalsize\bfseries}{\thesubsection}{0.6em}{}

\pagestyle{fancy}
\fancyhf{}
\fancyhead[L]{\small\leftmark}
\fancyfoot[C]{\small\thepage}
\renewcommand{\headrulewidth}{0.3pt}
\renewcommand{\chaptermark}[1]{\markboth{\thechapter.\ #1}{}}
\fancypagestyle{plain}{\fancyhf{}\fancyfoot[C]{\small\thepage}\renewcommand{\headrulewidth}{0pt}}

% ---- ambienti ----
\newtcolorbox{formula}[1][]{enhanced,breakable,colback=cGrigio,colframe=cBlu,boxrule=0.6pt,
  left=6pt,right=6pt,top=4pt,bottom=4pt,title={#1},fonttitle=\bfseries\small,
  coltitle=cBlu,attach boxed title to top left={yshift=-2mm,xshift=3mm},
  boxed title style={colback=white,colframe=cBlu,boxrule=0.4pt,sharp corners},
  sharp corners}
\newtcolorbox{attenzione}[1][Attenzione]{enhanced,breakable,colback=red!4,colframe=cRosso,boxrule=0.6pt,
  left=6pt,right=6pt,top=3pt,bottom=3pt,title={#1},fonttitle=\bfseries\small,coltitle=cRosso,
  attach boxed title to top left={yshift=-2mm,xshift=3mm},
  boxed title style={colback=white,colframe=cRosso,boxrule=0.4pt,sharp corners},sharp corners}
\newtcolorbox{nota}[1][Tip]{enhanced,breakable,colback=green!4,colframe=cVerde,boxrule=0.6pt,
  left=6pt,right=6pt,top=3pt,bottom=3pt,title={#1},fonttitle=\bfseries\small,coltitle=cVerde,
  attach boxed title to top left={yshift=-2mm,xshift=3mm},
  boxed title style={colback=white,colframe=cVerde,boxrule=0.4pt,sharp corners},sharp corners}
\newtcolorbox{procedura}[1][Procedura]{enhanced,breakable,colback=white,colframe=cArancio,boxrule=0.6pt,
  left=6pt,right=6pt,top=3pt,bottom=3pt,title={#1},fonttitle=\bfseries\small,coltitle=cArancio,
  attach boxed title to top left={yshift=-2mm,xshift=3mm},
  boxed title style={colback=white,colframe=cArancio,boxrule=0.4pt,sharp corners},sharp corners}
\newtcolorbox{esempio}[1]{enhanced,breakable,colback=white,colframe=black!55,boxrule=0.5pt,
  left=6pt,right=6pt,top=4pt,bottom=4pt,title={#1},fonttitle=\bfseries\small,
  colbacktitle=black!8,coltitle=black,sharp corners,toptitle=2pt,bottomtitle=2pt}

% ---- macro matematiche ----
\newcommand{\E}{\mathbb{E}}
\renewcommand{\Pr}{\mathrm{P}}
\newcommand{\Pb}{P_{\mathrm{b}}}
\newcommand{\la}{\lambda_{\mathrm{a}}}
\newcommand{\lr}{\lambda_{\mathrm{r}}}
\newcommand{\Nq}{N_{\mathrm{q}}}
\newcommand{\ie}{\emph{i.e.}\ }
\newcommand{\mio}{\textsuperscript{\textcolor{cArancio}{[m]}}}

% ---- catene di Markov in TikZ ----
\tikzset{
  st/.style={circle,draw,thick,minimum size=7.5mm,inner sep=0pt,font=\small},
  fw/.style={-{Stealth[length=2mm]},thick,bend left=18},
  bw/.style={-{Stealth[length=2mm]},thick,bend left=18},
  lp/.style={-{Stealth[length=2mm]},thick,red!70!black,loop above,looseness=6,min distance=5mm},
  lab/.style={font=\scriptsize,inner sep=1pt},
}

% catena nascita-morte: #1 ultimo stato, #2 lista i/etichetta archi su, #3 lista i/etichetta archi giu, #4 codice extra
\newcommand{\bdchain}[4]{%
\begin{tikzpicture}[baseline]
 \foreach \i in {0,...,#1} \node[st] (s\i) at (1.9*\i,0) {\i};
 \foreach \i/\l in {#2} {\pgfmathtruncatemacro{\j}{\i+1}\draw[fw] (s\i) to node[above,lab]{\l} (s\j);}
 \foreach \i/\l in {#3} {\pgfmathtruncatemacro{\j}{\i-1}\draw[bw] (s\i) to node[below,lab]{\l} (s\j);}
 #4
\end{tikzpicture}}
